\section{Task 2: Multi-Label Classification}
This task formulates glass-surface inspection as a \emph{multi-label} classification problem with four labels (no defect, chipped edge, scratch, stain). We train a convolutional neural network (ResNet backbone) with a sigmoid multi-label head and evaluate it using the \textbf{Micro F1-score} on the provided test set. With our current retraining run (2 epochs), the checkpoint achieves \textbf{Micro F1 = 0.8439} on the test set at a decision threshold of 0.50.

\subsection{Problem Formulation}
We define a label vector $\mathbf{y}\in\{0,1\}^4$ for each image with the following ordered labels:
\[
\mathbf{y} = [y_{\text{no\_defect}},\; y_{\text{chipped\_edge}},\; y_{\text{scratch}},\; y_{\text{stain}}].
\]

\subsubsection{Multi-label vs. multi-class}
This task is \emph{multi-label} (not multi-class): an image may contain multiple defect types simultaneously. Therefore, we model each label as an independent Bernoulli variable and use a sigmoid output for each label instead of a softmax over classes.

\subsubsection{Label Construction from Annotation Files}
Each image may have a corresponding \texttt{.txt} file that lists defect instances. We assign labels by the rules required in the task description:
\begin{itemize}
  \item If the image has \textbf{no defects}, then $y_{\text{no\_defect}}=1$ and the other three labels are 0.
  \item If the image has defects, then for each defect type present in the \texttt{.txt} file, we set the corresponding label to 1.
\end{itemize}

In our implementation, the dataset parser treats a missing \texttt{.txt} as a clean image (no defect). For images with an existing \texttt{.txt}, we mark each defect label as present if \emph{at least one} instance of that defect class appears in the file.

The annotation files use defect class IDs $\{0,1,2\}$, which we map to \texttt{chipped\_edge}, \texttt{scratch}, and \texttt{stain}, respectively. The \texttt{no\_defect} label is derived: it is set to 1 only when no defect ID appears (including the case of an empty file).

\subsubsection{Implementation (label parsing)}
The image-level label is computed by scanning the annotation file and marking whether each defect ID appears at least once. This is implemented in \texttt{Task2/dataset.py} (function \texttt{\_load\_label}). In words:
\begin{itemize}
  \item If the \texttt{.txt} file is missing, set \texttt{no\_defect=1}.
  \item Otherwise, set \texttt{chipped\_edge/scratch/stain} to 1 if the corresponding defect ID appears in any line.
  \item If no defect IDs appear (empty file), set \texttt{no\_defect=1}.
\end{itemize}

\subsection{Dataset and Preprocessing}
The dataset is organized as:\; \texttt{data/train/img}, \texttt{data/train/txt}, \texttt{data/test/img}, \texttt{data/test/txt}.

\subsubsection{What the \texttt{.txt} file means}
Each \texttt{.txt} file contains 0 or more lines of the form
\[\texttt{class\_id \; x\_center \; y\_center \; width \; height}\]
(the geometry fields are not needed for Task~2). We treat the image as having a defect label if \emph{any} line has a given defect class ID. This converts an instance-level annotation into an image-level multi-label target.

\subsubsection{Image Transformations}
Training uses standard resizing to $320\times 320$ and light augmentation (color jitter, flips, small rotations), followed by normalization. Validation/test use resizing and normalization only.

\subsubsection{Class Imbalance Handling}
Multi-label datasets are often highly imbalanced (e.g., far more \texttt{no\_defect} than rare defects). We address this using \texttt{BCEWithLogitsLoss} with a \texttt{pos\_weight} computed by scanning training labels, increasing the loss contribution of rare positives.

Concretely, for each label $k$, we compute
\[
\texttt{pos\_weight}_k = \frac{N_{\text{neg},k}}{N_{\text{pos},k}},
\]
where $N_{\text{pos},k}$ and $N_{\text{neg},k}$ are counted from training annotations. This encourages the optimizer not to ignore rare defect labels.

Figure~\ref{fig:prevalence} visualizes this imbalance: \texttt{no\_defect} is very frequent, while the defect labels are comparatively rare. This imbalance is the main reason that Micro F1 can appear high even when rare-label precision is limited.

\begin{figure}[t]
  \centering
  \includegraphics[width=\linewidth]{figures/label_prevalence.png}
  \caption{Label prevalence (positive rate) computed from annotation files.}
  \label{fig:prevalence}
\end{figure}

\subsection{Model}
We use a ResNet backbone and replace the final classification layer with a 4-unit linear head for multi-label prediction. The network outputs logits $\mathbf{z}\in\mathbb{R}^4$, converted to probabilities via a sigmoid:
\[
\hat{\mathbf{p}} = \sigma(\mathbf{z}).
\]
A label is predicted positive when $\hat{p}_k \ge t$ for threshold $t$.

\subsubsection{How the forward pass maps to labels}
Given an image $x$, the model predicts four independent probabilities
$\hat{p}_{\text{no\_defect}}, \hat{p}_{\text{chipped\_edge}}, \hat{p}_{\text{scratch}}, \hat{p}_{\text{stain}}$.
We use a sigmoid (not softmax) because multiple defect labels can be simultaneously true.

\subsubsection{Implementation (model)}
The model is implemented in \texttt{Task2/model.py} as \texttt{MultiLabelResNet}. We take a ResNet-34 backbone, remove its original classification layer, and add a small head that outputs 4 logits. During evaluation, we apply $\sigma(\cdot)$ and threshold the probabilities to obtain a 0/1 vector.

\subsubsection{Training Setup (Best Checkpoint)}
The checkpoint metadata indicates the following configuration:
\begin{itemize}
  \item Backbone: ResNet-34 with ImageNet initialization (pretrained)
  \item Image size: $320\times 320$
  \item Optimizer: AdamW, learning rate $3\times 10^{-4}$, weight decay $10^{-4}$
  \item Epochs: 2 (current run), cosine annealing learning-rate schedule
  \item Mixed precision (AMP): enabled (if CUDA available)
  \item Validation split: 10\% of training data
\end{itemize}

\begin{figure}[t]
  \centering
  \includegraphics[width=\linewidth]{figures/training_curves.png}
  \caption{Training curves recorded during training (epoch-level train/val loss, validation Micro F1, and learning rate).}
  \label{fig:training-curves}
\end{figure}

\subsubsection{Training log (per epoch)}
We record epoch-level metrics to \texttt{Task2/report/train\_log.json} and \texttt{Task2/report/train\_log.csv}. Table~\ref{tab:task2-trainlog} summarizes the current run.

\begin{table}[t]
\centering
\small
\begin{tabular}{rrrrr}
\toprule
Epoch & Train loss & Val loss & Val Micro F1 & Best $t$ (val)\\
\midrule
1 & 0.6368 & 0.6095 & 0.8242 & 0.70\\
2 & 0.4875 & 0.4822 & 0.8554 & 0.50\\
\bottomrule
\end{tabular}
\caption{Epoch-level training/validation metrics for the current run.}
\label{tab:task2-trainlog}
\end{table}

\subsubsection{What is optimized (loss function)}
We optimize a weighted binary cross-entropy over four labels. Using logits $z_k$ and targets $y_k\in\{0,1\}$, the per-label term is
\[
\ell_k = -\,w_k\,y_k\log\sigma(z_k)\; -\;(1-y_k)\log(1-\sigma(z_k)),
\]
where $w_k$ is the \texttt{pos\_weight} for label $k$. The final loss is the mean over labels and batch.

\subsubsection{Implementation (training loop and metric computation)}
Training is implemented in \texttt{Task2/train.py}. Each epoch performs:
\begin{itemize}
  \item One training pass that minimizes weighted BCE with logits.
  \item A validation pass that computes \textbf{Micro F1} using \texttt{Task2/utils.py} (function \texttt{evaluate}).
  \item A threshold search on the validation set via \texttt{threshold\_search} (tries candidate thresholds and selects the one with the best Micro F1).
\end{itemize}

To make the experiment reportable, \texttt{Task2/train.py} records epoch-level metrics to \texttt{train\_log.json/csv} and updates \texttt{training\_curves.png} after every epoch.

\subsection{Threshold Selection}
Since sigmoid outputs are continuous, we must choose a decision threshold $t$. We select $t$ by searching over candidate values on the validation split (0.20 to 0.80 with step 0.05) and storing the threshold that maximizes Micro F1.

In the current run, the best validation threshold changes across epochs (Table~\ref{tab:task2-trainlog}), illustrating that thresholding is part of model selection in multi-label settings.

In addition, we also plot Micro F1 as a function of threshold on the \emph{test set} to visualize sensitivity to $t$ (this is diagnostic only; the reported test score still uses the stored threshold from validation).

\begin{figure}[t]
  \centering
  \includegraphics[width=0.92\linewidth]{figures/threshold_scan_test.png}
  \caption{Micro F1 vs threshold on the test set (diagnostic).}
  \label{fig:threshold-scan-test}
\end{figure}

\subsection{Evaluation Metric: Micro F1}
We report \textbf{Micro F1-score} on the test set, as required. For multi-label classification, Micro F1 aggregates true positives, false positives, and false negatives over \emph{all} labels before computing F1, emphasizing overall instance-level correctness.

For completeness, we also report the macro F1 (unweighted mean over labels) to highlight performance on rare defect labels. Macro F1 is typically lower under heavy imbalance.

\subsection{Results on the Test Set}
Using the saved checkpoint and the stored threshold $t=0.50$, we obtain:
\[
\text{Micro F1 (test)} = 0.8439.
\]

We also compute the test-set weighted BCE loss (for reference):
\[
\text{BCE loss (test)} = 0.2423.
\]

\input{table/metrics_table.tex}

\paragraph{Key observations from Table~\ref{tab:task2-metrics}.}
\begin{itemize}
  \item \textbf{Imbalance dominates Micro F1.} The \texttt{no\_defect} label has very large support (9501 positives), so its performance strongly influences Micro F1.
  \item \textbf{The model prioritizes recall on rare defects.} For \texttt{chipped\_edge}, \texttt{scratch}, and \texttt{stain}, recall is high, but precision is much lower. This indicates many false positives for rare labels.
  \item \textbf{Macro F1 is much lower than Micro F1.} This is expected under heavy imbalance: macro F1 weights rare labels equally and therefore highlights their difficulty.
\end{itemize}

\paragraph{Why precision is low for rare labels.} There are two main factors. (1) Positive samples for rare defects are scarce compared with clean images, so the classifier may over-trigger defect labels. (2) Defect appearances can be subtle and share texture statistics with normal glass patterns, increasing confusion. Using \texttt{pos\_weight} improves recall but can also increase false positives; threshold tuning is therefore important.

\paragraph{Average Precision (AP) interpretation.} Figure~\ref{fig:pr-curves} shows Precision--Recall curves, which summarize performance across all thresholds. The AP values (area under the PR curve) further highlight the difficulty of rare labels: \texttt{no\_defect} is near-perfect, while \texttt{chipped\_edge} is substantially harder, suggesting that the raw probability ranking for this class is not very separable.

\begin{figure}[t]
  \centering
  \includegraphics[width=0.85\linewidth]{figures/per_label_f1.png}
  \caption{Per-label F1-score on the test set at threshold $t=0.50$.}
  \label{fig:per-label-f1}
\end{figure}

\begin{figure}[t]
  \centering
  \includegraphics[width=0.90\linewidth]{figures/pr_curves_test.png}
  \caption{Precision--Recall curves on the test set (includes micro-average).}
  \label{fig:pr-curves}
\end{figure}

\subsection{Discussion}
\paragraph{Strengths.} The model achieves strong Micro F1, driven largely by high performance on \texttt{no\_defect} and competitive recall on rarer defect classes. Threshold tuning provides a principled way to trade precision and recall.

\paragraph{Weaknesses.} Per-label precision for rare classes can be lower due to class imbalance and overlapping visual patterns. In multi-label settings, false positives on rare labels can meaningfully impact downstream inspection workflows.

In this experiment, the gap between Micro F1 and Macro F1 is consistent with an imbalanced dataset: the model is reliable on the dominant clean class, while defect classes (especially visually subtle ones) show a precision--recall trade-off.

\paragraph{Effect of training length.} This report reflects a short retraining run (2 epochs). Longer training typically improves the representation and calibration of rare labels, which can increase both Macro F1 and Micro F1. Therefore, the defect-label precision/recall trade-off observed here should be interpreted as an early-stage model rather than a fully converged one.

\paragraph{Potential Improvements.}
\begin{itemize}
  \item Use per-class thresholds (rather than a single global $t$) to better calibrate rare classes.
  \item Add stronger augmentations (e.g., RandAugment) and/or class-aware sampling.
  \item Use focal loss or asymmetric loss tailored for multi-label imbalance.
  \item Add probability calibration (temperature scaling) to improve decision reliability.
\end{itemize}

\subsection{Reproducibility}
To reproduce the reported Micro F1 on the test set:
\begin{itemize}
  \item Run evaluation: \texttt{python3 Task2/eval.py --data\_path data --checkpoint Task2/best\_task2\_retrain.pth --img\_size 320}
  \item Regenerate report assets (table/plots): \texttt{python3 Task2/report/compute\_task2\_metrics.py --data\_path data --checkpoint Task2/best\_task2\_retrain.pth --img\_size 320}
  \item Re-train and record curves: \texttt{python3 Task2/train.py --data\_path data --img\_size 320 --batch\_size 32 --epochs 2 --lr 0.0003 --weight\_decay 0.0001 --backbone resnet34 --pretrained --val\_split 0.1 --num\_workers 4 --seed 42 --amp --log\_interval 20 --plot\_curves --log\_dir Task2/report --save\_path Task2/best\_task2\_retrain.pth}
\end{itemize}

\subsubsection{Where it is implemented in code}
The pipeline is implemented as follows:
\begin{itemize}
  \item Dataset + label parsing: \texttt{Task2/dataset.py} (function \texttt{\_load\_label})
  \item Model definition: \texttt{Task2/model.py} (class \texttt{MultiLabelResNet})
  \item Training loop + threshold search: \texttt{Task2/train.py} (uses \texttt{threshold\_search} and \texttt{save\_checkpoint})
  \item Metric computation (Micro F1): \texttt{Task2/utils.py} (function \texttt{evaluate})
  \item Test evaluation entry point: \texttt{Task2/eval.py}
\end{itemize}